% from ega4-8-fr.tex
\oldpage[IV-3]{13}
\item[\emph{b)}] Il existe $\lambda\geq\alpha$ tel que
$u_\lambda^{-1}(E_\lambda)\subset u_\lambda^{-1}(F_\lambda)$ (et alors on a
$u_\mu^{-1}(E_\mu)\subset u_\mu^{-1}(F_\mu)$ pour $\mu\geq\lambda$).
\item[\emph{c)}] Il existe $\lambda\geq\alpha$ tel que
$E_\lambda\subset F_\lambda$ (et alors on a $E_\mu\subset F_\mu$ pour
$\mu\geq\lambda$).
\end{enumerate}
\end{theorem}

Les remarques entre parenthèses dans \emph{b)} et \emph{c)} résultent de (ii).
Posons
\[
G_\lambda=E_\lambda\cap(S_\lambda-F_\lambda),\qquad
G=E\cap(S-F).
\]
Alors $G_\lambda$ est une partie pro-constructible de $S_\lambda$ (1.9.5,
(i)), et en vertu de (8.3.2.1) et de (ii), on a
\[
G_\mu\subset u_{\lambda\mu}^{-1}(G_\lambda)
\quad\text{pour }\lambda\leq\mu,
\qquad
G=\bigcap_\lambda u_\lambda^{-1}(G_\lambda).
\]
On est donc ramené à prouver le cas particulier de (8.3.2) correspondant à
$F_\lambda=\emptyset$ pour tout $\lambda$~:

\begin{corollary}[8.3.3]
\label{IV.8.3.3-fr}
Pour tout $\lambda$, soit $E_\lambda$ une partie pro-constructible de
$S_\lambda$ telle que, pour $\lambda\leq\mu$, on ait
$E_\mu\subset u_{\lambda\mu}^{-1}(E_\lambda)$. Supposons qu'il existe
$\alpha\in L$ tel que $S_\alpha$ soit quasi-compact. Alors les conditions
suivantes sont équivalentes~:
\begin{enumerate}
\item[\emph{a)}]
$E=\bigcap_\lambda u_\lambda^{-1}(E_\lambda)=\emptyset$.
\item[\emph{b)}] Il existe un $\lambda$ tel que
$u_\lambda^{-1}(E_\lambda)=\emptyset$ (et alors
$u_\mu^{-1}(E_\mu)=\emptyset$ pour $\mu\geq\lambda$).
\item[\emph{c)}] Il existe $\lambda$ tel que $E_\lambda=\emptyset$ (et alors
$E_\mu=\emptyset$ pour $\mu\geq\lambda$).
\end{enumerate}
\end{corollary}

Il est clair que \emph{c)} implique \emph{a)}. Prouvons que \emph{a)} entraîne
\emph{b)}~: puisque $S_\lambda$ est quasi-compact, il en est de même de $S$
(8.2.2)~; $E_\lambda$ étant pro-constructible, il en est de même de
$u_\lambda^{-1}(E_\lambda)$ (1.9.5, (vi))~; alors la famille filtrante
décroissante d'ensembles pro-constructibles $u_\lambda^{-1}(E_\lambda)$ a
une intersection vide, donc (1.9.9) l'un d'eux est vide.

Enfin, montrons que \emph{b)} entraîne \emph{c)}. Comme $S_\alpha$ est
quasi-compact et $L$ filtrant, on peut remplacer $S_\alpha$ par un ouvert
affine, donc supposer (8.2.2) que $S$ et les $S_\lambda$ pour
$\lambda\geq\alpha$ sont affines~; on a alors (1.9.2.1), pour
$\lambda\geq\alpha$,
\[
u_\lambda(S)=\bigcap_{\mu\geq\lambda}u_{\lambda\mu}(S_\mu),
\]
d'où
\[
E_\lambda\cap u_\lambda(S)=
\bigcap_{\mu\geq\lambda}
\bigl(E_\lambda\cap u_{\lambda\mu}(S_\mu)\bigr).
\]
Or, comme $u_\lambda$ et les $u_{\lambda\mu}$ sont quasi-compacts,
$u_\lambda(S)$ et $u_{\lambda\mu}(S_\mu)$ sont pro-constructibles dans
$S_\lambda$ (1.9.5, (vii)), donc les ensembles
$E_\lambda\cap u_{\lambda\mu}(S_\mu)$ pour $\mu\geq\lambda$ forment une
famille filtrante décroissante de parties pro-constructibles de $S_\lambda$.
Comme $S_\lambda$ est quasi-compact, l'hypothèse \emph{b)} entraîne que
l'intersection de cette famille est vide, donc (1.9.9) un des ensembles
$E_\lambda\cap u_{\lambda\mu}(S_\mu)$ est vide, donc
$E_\mu\subset u_{\lambda\mu}^{-1}(E_\lambda)$ est vide. C.Q.F.D.

\begin{corollary}[8.3.4]
\label{IV.8.3.4-fr}
Pour tout $\lambda$, soit $F_\lambda$ une partie ind-constructible de
$S_\lambda$ telle que pour $\lambda\leq\mu$ on ait
$u_{\lambda\mu}^{-1}(F_\lambda)\subset F_\mu$. Supposons qu'il existe
$\alpha\in L$ tel que $S_\alpha$ soit quasi-compact. Alors les conditions
suivantes sont équivalentes~:
\begin{enumerate}
\item[\emph{a)}] L'ensemble
$F=\bigcup_\lambda u_\lambda^{-1}(F_\lambda)$ est égal à $S$.
\item[\emph{b)}] Il existe un $\lambda$ tel que
$u_\lambda^{-1}(F_\lambda)=S$ (et alors
$u_\mu^{-1}(F_\mu)=S$ pour $\mu\geq\lambda$).
